\begin{statement}[C.8.5; Proved: R4 by explicit analytic estimates]
For every lattice $\Z+\tau\Z$ and every zero $z$ of its
Weierstrass function,
\[
 N(z,\tau)\le\sqrt2.
\]
Equality holds exactly for the modular orbit of $i$.
In the parameter of the family this says
$\max D_{\mathrm{part}}=-(3/\pi)\log6$, only at $t=1$.
\end{statement}
\begin{proof}
Reduce to the standard modular region. The portion $T\ge2$
is covered by Statement C.8.1. Suppose $T\le2$ and put
$r=|t-1|$. For $r\ge9/10$, the step $n=2$ in
\eqref{eq:ca-ladder} and $N(2P)<17/10$ give strict inequality,
because
\[
 \frac{27(17/10)^4}{256}<\frac9{10}.
\]
There is no $2$-torsion point of the section on the smooth base.

If $r\le9/10$, then $T<7/5$. Indeed, $T\ge7/5$ would give
$|q|<1/6500$ and
\[
 |E_4|>24/25,\qquad |\Delta_0|<1/6400,
 \qquad |t|>\frac{6400(24/25)^3}{1728}>19/10,
\]
contrary to $|t|\le19/10$. The estimates follow from
\[
 \sum_{n\ge1}\sigma_3(n)r^n
 \le\frac{r(1+4r+r^2)}{(1-r)^5},\qquad
 |\Delta_0|\le r\exp\frac{24r}{1-r}.
\]
For $1/28\le r\le9/10$, the $n=4$ step and
$N(4P)<3/2$ now apply. The function
$(1-5r/27)r^{1/4}$ is increasing up to $r=27/25$.
At its left endpoint its fourth power is
\[
 \frac1{28}\left(\frac{751}{756}\right)^4
 >\frac{2187}{65536}
 =\left(\frac{3^{3/4}}8\frac32\right)^4.
\]
This proves strict inequality on the whole annulus. The lower
bound for $A_2$ is positive there, so no division by a zero has
occurred.

If $r\le1/28$, then $T<11/10$. For otherwise $|q|<1/1000$,
and the Lambert series with
\[
 \sum_{n\ge1}\sigma_5(n)r^n
 \le\frac{r(1+26r+66r^2+26r^3+r^4)}{(1-r)^7}
\]
gives $|E_6|>9/20$, $|\Delta_0|<1/900$, and hence
$|t-1|>900(9/20)^2/1728>1/28$.
On $1/500\le r\le1/28$, the $n=8$ step,
$N(8P)<143/100$, and $|A_3|\ge1-10r$ prove the claim.
Indeed $(1-10r)r^{1/4}$ has a single critical point in this
interval, a maximum at $r=1/50$, so its minimum is at an endpoint.
At both endpoints its fourth power exceeds $9/5000$, whereas
\[
 \left(\frac{3^{3/4}}{16}\frac{143}{100}\right)^4
 <\frac7{4000}<\frac9{5000}.
\]
Again the displayed positive bound excludes a division-polynomial
zero. The remaining disc $r\le1/500$ is covered by Statement
C.8.3. The smooth limiting fibre at $t=0$ has
$N=3^{1/4}<\sqrt2$; the formula also extends there by
continuity, so the singular rational model causes no omitted
modulus. The only equality occurs at $t=1$.
\end{proof}
